\documentclass[11pt]{article}
\usepackage[a4paper,margin=25mm]{geometry}
\usepackage[T1]{fontenc}
\usepackage{lmodern}
\usepackage{amsmath,amssymb,amsthm,booktabs}
\usepackage[hidelinks]{hyperref}
\hypersetup{pdftitle={Essential singleton Seidel switches and a construction from arbitrary graphs},pdfauthor={Haobo Ma and Wenlin Zhang}}
\newtheorem{theorem}{Theorem}
\newtheorem{proposition}[theorem]{Proposition}
\newtheorem{lemma}[theorem]{Lemma}
\newtheorem{corollary}[theorem]{Corollary}
\newcommand{\ISS}{\mathrm{ISS}}
\title{Essential singleton Seidel switches and a construction from arbitrary graphs}
\author{Haobo Ma\\{\small ChronoAI Pte Ltd; The Omega Institute}
\and Wenlin Zhang\\{\small National University of Singapore; The Omega Institute}}
\date{}
\begin{document}
\maketitle
\begin{abstract}
A graph is covered if every vertex belongs to an identity Seidel switch
consisting of a singleton or the endpoints of an edge. We prove that, in a
covered graph of order $2h+1$, $h\geq1$, with degrees in $\{h,h+1\}$, every vertex
requiring a singleton switch is the unique singleton-switch vertex and is
fixed by every isomorphism from its switch to the original graph. From each
graph $F$ on at least two vertices without isolated vertices we construct a covered graph on
$2|V(F)|+1$ vertices with exactly one essential singleton, and determine all
its singleton and edge switches. Nonisomorphic base graphs give
nonisomorphic constructed graphs; the construction starts with the house.
An exhaustive check of orders three through eight finds seven covered graphs
with essential singletons; each has a unique singleton-ISS vertex, fixed by
every switching isomorphism. The same behavior holds in the sixteen
essential-singleton cases at order nine.
\end{abstract}

\section{The fixing question}
Let $G=(V,E)$ be a finite simple graph. For $S\subseteq V$, its Seidel
switch $G^S$ toggles precisely the adjacencies between $S$ and $V\setminus S$.
Following Gervacio~\cite{gervacio}, call $S$ an \emph{identity Seidel switch}
(ISS) if $G^S\cong G$. A singleton ISS is a vertex-ISS; an edge-ISS consists
of two adjacent vertices. A graph is \emph{covered} if every vertex lies
in a singleton ISS or an edge-ISS. Set
\[
 W(G)=\{v:\{v\}\text{ is an ISS}\},\qquad
 E(P(G))=\{xy\in E(G):\{x,y\}\text{ is an ISS}\},
\]
and let $U(G)$ be the set of isolated vertices of $P(G)$. In a covered
graph, $U(G)\subseteq W(G)$, and these vertices require singleton switches.
All switching isomorphisms below have direction $G^S\longrightarrow G$.

Gervacio~\cite{communication} proposed studying $U(G)$ and asked whether
in a covered graph every $v\in U(G)$ admits a switching isomorphism
fixing $v$, suggesting
$|U(G)|=1$ as a first case and the house as the initial example. He also
gave the exact criterion that such an isomorphism exists if and only if
some automorphism of $G-v$ exchanges $N_G(v)$ and
$V\setminus(\{v\}\cup N_G(v))$. Indeed, restriction gives the required
automorphism, and extension by $v\mapsto v$ gives the converse.

We answer the question affirmatively when the degrees lie in
$\{h,h+1\}$, where $|V|=2h+1$, and construct a family with $|U(G)|=1$
from every graph on at least two vertices without isolated vertices.
The question for arbitrary
covered graphs, including arbitrary covered graphs with $|U(G)|=1$,
remains open. Our arguments use switching covariance and exact degree
changes; they do not assume that the ISS subsets form a group.

\section{Tracking singleton switches}
For a permutation $p$ of $V$, write $pG$ for the relabelled graph, so that
$p(G^S)=(pG)^{p(S)}$. Also $(G^S)^T=G^{S\triangle T}$.
These identities concern labelled graphs and hold for all subsets.

\begin{lemma}\label{transport}
Suppose $v\in W(G)$ and $p:G^{\{v\}}\to G$ is an isomorphism. Then
$p(v)\in W(G)$. In particular, if $W(G)=\{v\}$, every such $p$ fixes $v$.
Moreover,
\begin{equation}\label{localW}
 W(G^{\{v\}})=\{v\}\cup
 \{x\ne v:\{v,x\}\text{ is an ISS of }G\}.
\end{equation}
For every such $v$ and $p$, with $k=|W(G)|$,
\begin{equation}\label{partnercount}
 d_{P(G)}(v)=k-1-d_{G[W(G)]}(p(v)).
\end{equation}
In particular, if $v\in U(G)$, then $p(v)$ is adjacent to every other
vertex of $W(G)$. Consequently, if $W(G)$ is independent and $v\in W(G)\cap U(G)$, then
$W(G)=\{v\}$ and every switching isomorphism at $v$ fixes $v$.
\end{lemma}
\begin{proof}
Covariance gives $(pG)^{\{p(v)\}}=G$, hence
$pG=G^{\{p(v)\}}$. Thus $G^{\{p(v)\}}\cong G$, proving the first claim.
Since $G^{\{v\}}\cong G$, for $x\ne v$ the singleton $x$ is an ISS
of $G^{\{v\}}$ exactly when $G^{\{v,x\}}\cong G$, which proves
\eqref{localW}; switching $v$ a second time returns $G$.

If $v\in U(G)$, every $x\ne v$ on the right of \eqref{localW} is a
nonneighbor of $v$ in $G$, and therefore a neighbor in $G^{\{v\}}$.
The isomorphism $p$ takes $W(G^{\{v\}})$ onto $W(G)$, so $p(v)$ is
adjacent to every other vertex of $W(G)$. In general, the edge-ISS
partners of $v$ in $G$ are precisely its nonneighbors in
$W(G^{\{v\}})\setminus\{v\}$, which proves \eqref{partnercount}.
An independent graph with a
vertex adjacent to all its other vertices has just one vertex.
\end{proof}
The distinction between $|U(G)|=1$ and $|W(G)|=1$ matters: the lemma
proves the fixing conclusion from the latter hypothesis. It does not
identify the two conditions in arbitrary covered graphs.
If $G[W(G)]$ is complete, \eqref{partnercount} gives $U(G)=W(G)$
for a covered graph. Hence a covered graph with $|U(G)|=1$ and
$|W(G)|>1$ must have $|W(G)|\geq3$, and its induced singleton graph
can be neither complete nor independent.

Switching a singleton changes the edge count by
$|V|-1-2d_G(v)$. Switching an adjacent pair $xy$ changes it by
$2(|V|-d_G(x)-d_G(y))$. Thus
\begin{equation}\label{degrees}
 v\in W(G)\ \Longrightarrow\ 2d_G(v)=|V|-1,
 \qquad xy\in E(P(G))\ \Longrightarrow\ d_G(x)+d_G(y)=|V|.
\end{equation}
In particular a covered graph with $U(G)\ne\varnothing$ has odd order.

\begin{theorem}\label{twolevels}
Let $G$ have order $2h+1$, where $h\geq1$, and suppose every degree is
$h$ or $h+1$. If $v\in W(G)\cap U(G)$, then $W(G)=\{v\}$ and every
isomorphism $G^{\{v\}}\to G$ fixes $v$. Thus Gervacio's fixing question
has an affirmative answer for every covered graph with these degrees.
Furthermore, in any such graph an edge $xy$ is an edge-ISS if and only if
every other vertex is adjacent to exactly one of $x,y$.
\end{theorem}
\begin{proof}
Take $v\in W(G)$. By \eqref{degrees}, $d_G(v)=h$. A neighbor of $v$
of degree $h$ would have degree $h-1$ after switching, and a nonneighbor
of degree $h+1$ would have degree $h+2$. Both are impossible in a graph
isomorphic to $G$. It follows that
\begin{equation}\label{neighbors}
 N_G(v)=\{x:d_G(x)=h+1\}.
\end{equation}
Every member of $W(G)$ has degree $h$, so $W(G)$ is independent.
Lemma~\ref{transport} gives the fixing assertion.

Now let $xy$ be an edge-ISS. For $z\notin\{x,y\}$, put
$c_z=|N_G(z)\cap\{x,y\}|$. Its new degree is
$d_G(z)+2-2c_z$. If $c_z=0$ or $2$, this falls outside $\{h,h+1\}$.
Thus $c_z=1$ for every such $z$. Conversely, if every outside vertex
sees exactly one endpoint, the transposition $(x\ y)$, fixing all
other vertices, is an isomorphism $G^{\{x,y\}}\to G$.
\end{proof}

\section{A family with one essential singleton}
Let $F$ be any simple graph on $\{1,\ldots,m\}$, $m\geq2$. Construct
$C(F)$ on
\[
 \{v\}\cup A\cup B,\qquad
 A=\{a_1,\ldots,a_m\},\quad B=\{b_1,\ldots,b_m\},
\]
as follows. The vertex $v$ is adjacent to all of $A$ and none of $B$.
Both $A$ and $B$ induce copies of $F$. Between the copies, put
\begin{equation}\label{cross}
 a_i b_j\in E(C(F))\quad\Longleftrightarrow\quad
 i=j\text{ or }ij\notin E(F).
\end{equation}
For distinct indices the second alternative has its usual meaning;
in particular every $a_i b_i$ is an edge.

\begin{theorem}\label{construction}
For every $F$ as above, $C(F)$ is covered and has degrees
\[
 d(v)=m,\qquad d(a_i)=m+1,\qquad d(b_i)=m.
\]
Its singleton ISS vertices are exactly
\begin{equation}\label{familyW}
 W(C(F))=\{v\}\cup\{b_i:d_F(i)=0\}.
\end{equation}
Its edge-ISS are exactly
\begin{equation}\label{familyP}
 \{a_i b_j:N_F(i)=N_F(j)\}
 \ \cup\ \{v a_i:d_F(i)=0\}.
\end{equation}
Consequently,
\[
 U(C(F))=
 \begin{cases}
 \{v\},&\text{if $F$ has no isolated vertices},\\
 \varnothing,&\text{otherwise}.
 \end{cases}
\]
If $F$ has no isolated vertices, every switching isomorphism at $v$
fixes $v$. Nonisomorphic graphs $F$ without isolated vertices yield
nonisomorphic graphs $C(F)$.
\end{theorem}
\begin{proof}
The degree formulas follow from
$d_F(i)+(m-d_F(i))=m$. The permutation
\[
 \tau(v)=v,\qquad \tau(a_i)=b_i,\qquad \tau(b_i)=a_i
\]
is an automorphism of $C(F)-v$ exchanging $A$ and $B$: the internal
copies agree, and the condition in \eqref{cross} is symmetric in $i,j$.
Extending it by fixing $v$ certifies the singleton switch at $v$.
For each $i$, the vertices $a_i,b_i$ are complementary outside their
pair, including at $v$. Their transposition certifies an edge-ISS.
These switches cover every vertex.

By \eqref{degrees}, no $a_i$ is a singleton ISS. If $b_i$ is one,
\eqref{neighbors} requires $N(b_i)=A$, which holds precisely when $i$
is isolated in $F$. Conversely, for such an $i$, the vertices $v,b_i$
have identical open neighborhoods $A$, so their transposition is an
automorphism of $C(F)$. Relabelling the singleton ISS at $v$ therefore
gives one at $b_i$. This proves \eqref{familyW}.

By Theorem~\ref{twolevels}, edge-ISS are exactly the adjacent pairs
complementary outside their pair. The degree sum in \eqref{degrees}
restricts them to $a_i b_j$ or $v a_i$. For an adjacent pair $a_i b_j$ with $i\ne j$, we have
$ij\notin E(F)$. For each $k\notin\{i,j\}$, complementarity at
$a_k$ (equivalently, at $b_k$) requires
$\mathbf{1}_{ik\in E(F)}=\mathbf{1}_{jk\in E(F)}$. At $v,a_j,b_i$
complementarity then holds automatically. Thus the two adjacency rows
agree, and conversely their equality gives complementarity everywhere.
For $i=j$ complementarity was already checked. This gives
$N_F(i)=N_F(j)$; if $i\ne j$
it also implies $ij\notin E(F)$, so the proposed pair is adjacent.
For $v a_i$, complementarity requires $a_i$ to have no neighbor in
$A\setminus\{a_i\}$ and all of $B$ as neighbors, which is exactly
$d_F(i)=0$. This proves \eqref{familyP}.

Every $a_i,b_i$ lies in the edge-ISS $a_i b_i$. Formula
\eqref{familyP} decides whether $v$ has an edge-ISS partner, giving
the assertion about $U$. If $F$ has no isolated vertices,
\eqref{familyW} and Lemma~\ref{transport} show that every switching
isomorphism at $v$ fixes it. Finally, any isomorphism between two such
constructed graphs must send their unique singleton-ISS vertices to
each other, hence their neighborhoods to each other. Those
neighborhoods induce the respective base graphs $F$.
\end{proof}

Taking $F=K_2$ gives the house: $C(F)-v$ is a four-cycle, and $v$ is
joined to the endpoints of one of its edges. Taking $F=K_m$ proves
existence with $|U|=1$ at every odd order at least five. More generally,
if $a_m$ denotes the number of isomorphism classes of $m$-vertex graphs
without isolated vertices, the theorem gives at least $a_m$ covered
isomorphism classes of order $2m+1$ with $|U|=1$.

The degree-sum test does not explain essentiality in this family.
When $F$ has no isolated vertices, each candidate $v a_i$ passes that
test. Choose $j\in N_F(i)$. The vertex $a_j$ is a common neighbor of
$v,a_i$, so switching that pair changes its degree from $m+1$ to
$m-1$, which is absent from $C(F)$. This explicitly rejects every
possible edge partner of $v$.

\section{Finite catalogues and the remaining question}
We recomputed all singleton and adjacent-pair switches for orders three
through eight, using the NetworkX distribution~\cite{networkxatlas} of the
Read--Wilson atlas~\cite{atlas} through order seven and McKay's
catalogue~\cite{mckay} at order eight. These calculations give:
\begin{center}
\begin{tabular}{rrrrr}
\toprule
Order & Input classes & Covered & $U=\varnothing$ & $|U|=1$\\
\midrule
3 & 4 & 1 & 1 & 0\\
4 & 11 & 3 & 3 & 0\\
5 & 34 & 3 & 2 & 1\\
6 & 156 & 11 & 11 & 0\\
7 & 1044 & 10 & 4 & 6\\
8 & 12346 & 35 & 35 & 0\\
\bottomrule
\end{tabular}
\end{center}
No covered graph in this range has $|U|\geq2$. In all seven with
$U\ne\varnothing$, the unique essential vertex is also the unique
singleton-ISS vertex. Thus Lemma~\ref{transport} explains why every
switching isomorphism there fixes it. The following table gives graph6
labels, the essential vertex (numbered from zero), and one certificate.
The tuple lists the images of vertices in order.
\begin{center}
\small
\begin{tabular}{lllr}
\toprule
Graph6 & Vertex & Fixing permutation & All isomorphisms\\
\midrule
\texttt{Dlc} & 4 & $(1,0,3,2,4)$ & 2\\
\texttt{FcBzo} & 0 & $(0,2,1,4,3,6,5)$ & 4\\
\texttt{FfxbO} & 0 & $(0,2,1,5,6,3,4)$ & 1\\
\texttt{FllHo} & 0 & $(0,5,4,6,2,1,3)$ & 1\\
\texttt{FBfnO} & 0 & $(0,5,6,4,3,1,2)$ & 4\\
\texttt{F\textasciigrave urg} & 2 & $(3,5,2,1,6,4,0)$ & 6\\
\texttt{FHN]w} & 1 & $(2,1,0,5,6,3,4)$ & 4\\
\bottomrule
\end{tabular}
\end{center}
Each successful switch was verified on every vertex pair. Negative
isomorphism decisions use exhaustive degree-compatible bijection search.
For the seven essential cases, a second check tested all $n!$ permutations
and confirmed the displayed counts and the fixing property. Catalogue
completeness uses the cited input catalogues. As a separate check of
Theorem~\ref{construction}, all $1098$ labelled base graphs of orders two
through five were constructed and their complete singleton and edge-ISS
lists compared with \eqref{familyW}--\eqref{familyP}.

An exhaustive supplementary check of McKay's $274668$ graphs of order
nine gives $33$ covered graphs: $17$ with $U=\varnothing$ and $16$ with
$U=W=\{v\}$. In all sixteen, every switching isomorphism at $v$ fixes
$v$. No covered graph of order nine has $|U|\geq2$. At the small orders
excluded from the tables, the one-vertex graph has $U=W=V$ and its
switch is the identity; the only covered graph of order two is $K_2$,
with $U=\varnothing$.

\newpage
These finite results leave two structural questions. Does every covered
graph with $|U|=1$ have a unique singleton-ISS vertex, or at least a fixing
switching isomorphism? Can a covered graph have $|U|\geq2$, and if so,
which graphs can occur as $G[U]$? The construction answers the existence
problem for $|U|=1$ at every odd order at least five, while the
fixing question beyond Theorem~\ref{twolevels} remains unresolved.

\begin{thebibliography}{9}
\bibitem{gervacio} S. V. Gervacio, \emph{On identity Seidel switches},
arXiv:2601.04530v1, 2026. \url{https://arxiv.org/abs/2601.04530v1}.
\bibitem{communication} S. V. Gervacio, private communication,
12 October 2026.
\bibitem{atlas} R. C. Read and R. J. Wilson, \emph{An Atlas of Graphs},
Oxford University Press, 1998.
\bibitem{networkxatlas} NetworkX developers, \emph{Graph atlas},
\texttt{graph\_atlas\_g} documentation; atlas data version 3.4.2.
\url{https://networkx.org/documentation/stable/reference/generated/networkx.generators.atlas.graph_atlas_g.html}.
\bibitem{mckay} B. D. McKay, \emph{Graphs}, small graph data catalogue.
\url{https://users.cecs.anu.edu.au/~bdm/data/graphs.html}.
\end{thebibliography}
\end{document}
